\documentclass[10pt]{article}

\usepackage[utf8]{inputenc}

\usepackage[T1]{fontenc}

\usepackage{amsmath}

\usepackage{amsfonts}

\usepackage{amssymb}

\usepackage[version=4]{mhchem}

\usepackage{stmaryrd}



\begin{document}

(K. Weierstrass). Вместо $\vartheta_{0}(v)$ часто пипут $\vartheta_{4}(v)$, имеются и другие системы обозначений. Сам $\mathrm{K}$. Якоби применял обозначения: $\quad \Theta(v)=\vartheta_{0}(v / 2 K), \quad H(v)=\vartheta_{1}(v / 2 K)$, $H_{1}(v)=\vartheta_{2}(v / 2 K), \quad \Theta_{1}(v)=\vartheta_{3}(v / 2 K), \quad$ где $K=\pi \vartheta_{3}^{2}(0) / 2$.



Все тета-функции Якоби представляют собой целые трансцендентные функции комплексного переменного $v$, причем $\vartheta_{1}(v)$ - нечетная функция, а остальные функции $\vartheta_{0}(v), \vartheta_{2}(v), \vartheta_{3}(v)$ четные.



Имеют место следующие с о о т н ш е н и я п ер и о д и ч н о с т и:



\[

\begin{aligned}

& \vartheta_{0}(v \pm 1)=\vartheta_{0}(v), \\

& \vartheta_{0}(v \pm \tau)=-e^{-i \pi \tau} \cdot e^{\mp 2 i \pi v} \cdot \vartheta_{0}(v) ; \\

& \vartheta_{1}(v \pm 1)=-\vartheta_{1}(v), \\

& \vartheta_{1}(v \pm \tau)=-e^{-i \pi \tau} \cdot e^{\mp 2 i \pi v} \cdot \vartheta_{1}(v) ; \\

& \vartheta_{2}(v \pm 1)=-\vartheta_{2}(v), \\

& \vartheta_{2}(v \pm \tau)=e^{-i \pi \tau} \cdot e^{\mp 2 i \pi v} \cdot \vartheta_{2}(v) ; \\

& \vartheta_{3}(v \pm 1)=\vartheta_{3}(v), \\

& \vartheta_{3}(v \pm \tau)=e^{-i \pi \tau} \cdot e^{\mp 2 i \pi v} \cdot \vartheta_{3}(v),

\end{aligned}

\]



из к-рых вытекает, что тета-функции являются эллиптич. функциями III рода по Эрмиту.



Различные тета-функции связаны между собой ф о рм улами п р е о б р а з о в а ия:



\[

\begin{aligned}

& \vartheta_{0}\left(v \pm \frac{1}{2}\right)=\vartheta_{3}(v) \\

& \vartheta_{0}\left(v \pm \frac{1}{2} \tau\right)= \pm i e^{-\frac{1}{4} i \pi \tau} \cdot e^{\mp i \pi v} \cdot \vartheta_{1}(v) \\

& \vartheta_{i}\left(v \pm \frac{1}{2}\right)= \pm \vartheta_{2}(v), \\

& \vartheta_{1}\left(v \pm \frac{1}{2} \tau\right)= \pm i e^{-\frac{1}{4} i \pi \tau} \cdot e^{\mp i \pi v} \cdot \vartheta_{0}(v) \\

& \vartheta_{2}\left(v \pm \frac{1}{2}\right)=\mp \vartheta_{1}(v), \\

& \vartheta_{2}\left(v \pm \frac{1}{2} \tau\right)=e^{-\frac{1}{4} i \pi \tau} \cdot e^{\mp i \pi v} \cdot \vartheta_{3}(v) \\

& \vartheta_{3}\left(v \pm \frac{1}{2}\right)=\vartheta_{0}(v), \\

& \vartheta_{3}\left(v \pm \frac{1}{2} \tau\right)=e^{-\frac{1}{4} i \pi \tau} \cdot e^{\mp i \pi v} \cdot \vartheta_{2}(v)

\end{aligned}

\]



Bсе четыре тета-функции удовлетворяют одному и тому жө дифференциальному уравнению



\[

\frac{\partial^{2} \vartheta}{\partial v^{2}}=4 i \pi \frac{\partial \vartheta}{\partial \tau}

\]



Важное значение имеют так наз. н у л е в ы е 3 н ачени и т е т а - ф у н к ц и й $\vartheta_{0}(0), \vartheta_{1}^{\prime}(0), \vartheta_{2}(0), \vartheta_{3}(0)$; при этом $\vartheta_{1}(0)=0$. Между ними имеются следующие соотношения:



\[

\vartheta_{1}^{\prime}(0)=\pi \vartheta_{0}(0) \vartheta_{2}(0) \vartheta_{3}(0), \quad \vartheta_{3}^{4}(0)=\vartheta_{0}^{4}(0)+\vartheta_{2}^{4}(0) .

\]



По отдельности



\[

\begin{gathered}

\vartheta_{0}(0)=H_{0} H_{3}^{2}, \quad \vartheta_{2}(0)=2 e^{\frac{1}{4} i \pi \tau} H_{0} H_{1}^{2}, \\

\vartheta_{3}(0)=H_{0} H_{2}^{2}, \quad \vartheta_{1}^{\prime}(0)=2 \pi e^{\frac{1}{4} i \pi \tau} H_{0}^{3}

\end{gathered}

\]



где



\[

\begin{aligned}

& H_{0}=\prod_{m=1}^{\infty}\left(1-e^{2 i \pi m \tau}\right), \\

& H_{1}=\prod_{m=1}^{\infty}\left(1+e^{2 i \pi m \tau}\right), \\

& H_{2}=\prod_{m=1}^{\infty}\left(1+e^{(2 m-1) i \pi \tau}\right), \\

& H_{3}=\prod_{m=1}^{\infty}\left(1-e^{(2 m-1) i \pi \tau}\right), \\

& H_{1} H_{2} H_{3}=1

\end{aligned}

\]



Функция $\vartheta_{0}(v)$ имеет простые нули в точках $m+$ $+\left(n-\frac{1}{2}\right) \tau ; \vartheta_{1}(v)$ - в точках $m+n \tau ; \vartheta_{2}(v)-$ в точках $m-\frac{1}{2}+n \tau ; \quad \vartheta_{3}(v)-\mathbf{\text { в }}$ точках $\quad m-\frac{1}{2}+\left(n-\frac{1}{2}\right) \tau ;$ $m, n=0, \pm 1, \ldots$



Из соотношений периодичности видно, что нек-рые отношения тета-функций будут эллиптич. Функциями в собственном смысле. О с н о вны е ә л л и п т ич е с ки е ф у нг ци и Я к о б и- $\operatorname{sn} u(c u н y c a m-$ плитуды), спи (косинус амплитуды) и $\operatorname{dn} u$ (дельта амплитудыс). Эти обозначения введены Х. Гудерманом (Ch. Gudermann, 1838). Названия происходят от старых обозначений, введенных самим $\mathrm{K}$. Якоби (sn $u=$ $=\sin \operatorname{am} u, \quad \operatorname{cn} u=\cos \operatorname{am} u, \operatorname{dn} u=\Delta \operatorname{am} u)$ и позднее вышедших из употребления.



Новое переменное $u$ связано с $v$ соотношением $u=$ $=v \cdot \pi \vartheta_{3}^{2}(0)$. Если $k=\vartheta_{2}^{2}(0) / \vartheta_{3}^{2}(0)-$ м о д у л ь э л ли пт и ч е с к и х ф у к ц и й, то Я. э. ф. следующим образом выражаются через тета-функции или посредством сходящихся в окрестности начала стешенных рядов:



$\operatorname{sn} u=\operatorname{sn}(u ; k)=\frac{\vartheta_{3}(0)}{\vartheta_{2}(0)} \frac{\vartheta_{1}(v)}{\vartheta_{0}(v)}=$



$=u-\left(1+k^{2}\right) \frac{u^{3}}{3 !}+\left(1+14 k^{2}+k^{4}\right) \frac{u^{5}}{5 !}-\ldots$,



$\operatorname{cn} u=\operatorname{cn}(u ; k)=\frac{\vartheta_{0}(0)}{\vartheta_{2}(0)} \frac{\vartheta_{2}(v)}{\vartheta_{0}(v)}=$



$=1-\frac{u^{2}}{2 !}+\left(1+4 k^{2}\right) \frac{u^{4}}{4 !}-\left(1+44 k^{2}+16 k^{4}\right) \frac{u^{6}}{6 !}+\ldots$,



$\operatorname{dn} u=\operatorname{dn}(u ; k)=\frac{\vartheta_{0}(0)}{\vartheta_{3}(0)} \frac{\vartheta_{3}(v)}{\vartheta_{0}(v)}=$



$=1-k^{2} \frac{u^{2}}{2 !}+k^{2}\left(4+k^{2}\right) \frac{u^{4}}{4 !}-k^{2}\left(16+44 k^{2}+k^{4}\right) \frac{u^{6}}{6 !}+\ldots$



Удобные обозначения для обратных величин и отношений были введены Дж. Глейшером (J. Glaisher, 1882):



\[

\begin{aligned}

\operatorname{ns} u=\frac{1}{\operatorname{sn} u}, & \text { nc } u=\frac{1}{\operatorname{cn} u}, \quad \text { nd } u=\frac{1}{\operatorname{dn} u}, \\

\operatorname{cs} u=\frac{\operatorname{cn} u}{\operatorname{sn} u}, & \operatorname{ds} u=\frac{\operatorname{dn} u}{\operatorname{sn} u}, \quad \operatorname{dc} u=\frac{\operatorname{dn} u}{\operatorname{cn} u}, \\

\operatorname{sc} u=\frac{\operatorname{sn} u}{\operatorname{cn} u}, & \text { sd } u=\frac{\operatorname{sn} u}{\operatorname{dn} u}, \quad \operatorname{cd} u=\frac{\operatorname{cn} u}{\operatorname{dn} u} .

\end{aligned}

\]



Я. ә. Ф. $\operatorname{sn} u, \operatorname{cn} u, \operatorname{dn} u$ являются эллиптич. функциями 2-го порядка с периодами: $4 K$ и $2 i K^{\prime}$ для $\operatorname{sn} u ; 4 K$ и $2\left(K+i K^{\prime}\right)$ для $\operatorname{cn} u ; 2 K$ и $4 i K^{\prime}$ для $\operatorname{dn} u$. Здесь $K=$ $=K(k)=\pi \vartheta_{3}^{2}(0) / 2, \quad K^{\prime}=K\left(k^{\prime}\right)=-i \tau K-$ значения полных эллиптич. интегралов I рода, $k^{\prime}=\sqrt{1-k^{2}}-$ д ополните льны й модуль эл липти ч ес к и х ф у нк ц и й. Я. ә. ф. имеют только простые полюсы, расположенные в точках $2 m K+(2 n+1) i K^{\prime}$;



Три Я. э. ф. связаны соотнотениями:



\[

\begin{gathered}

\operatorname{sn}^{2} u+\operatorname{cn}^{2} u=1, \\

\operatorname{dn}^{2} u+k^{2} \operatorname{sn}^{2} u=1, \\

\operatorname{dn}^{2} u-k^{2} \operatorname{cn}^{2} u=1-k^{2}=k^{\prime 2} \\

\frac{d}{d u} \operatorname{sn} u=\operatorname{cn} u \cdot \operatorname{dn} u \\

\frac{d}{d u} \operatorname{cn} u=-\operatorname{sn} u \cdot \operatorname{dn} u, \\

\frac{d}{d u} \operatorname{dn} u=-k^{2} \operatorname{sn} u \cdot \operatorname{cn} u

\end{gathered}

\]



и дифференциальными уравнениями



\[

\begin{gathered}

\left(\frac{d}{d u} \operatorname{sn} u\right)^{2}=\left(1-\operatorname{sn}^{2} u\right)\left(1-k^{2} \operatorname{sn}^{2} u\right), \\

\left(\frac{d}{d u} \operatorname{cn} u\right)^{2}=\left(1-\operatorname{cn}^{2} u\right)\left(k^{\prime 2}+k^{2} \operatorname{cn}^{2} u\right), \\

\left(\frac{d}{d u} \operatorname{dn} u\right)^{2}=\left(1-\operatorname{dn}^{2} u\right)\left(\operatorname{dn}^{2} u-k^{\prime 2}\right) .

\end{gathered}

\]





\end{document}